\documentclass[12pt]{article}
\usepackage[utf8]{inputenc}
\usepackage[T1]{fontenc}
\usepackage[margin=1in]{geometry}
\usepackage{setspace}
\usepackage{xcolor}
\usepackage[colorlinks=true, urlcolor=blue!60!black, linkcolor=blue!60!black]{hyperref}
\setlength{\parindent}{2em}
\setlength{\parskip}{1.25em}
\renewcommand{\baselinestretch}{1.0}
\begin{document}
\begin{center}
{\LARGE \textit{Employment Division v. Smith}, 494 U.S. 872 (1990)\par}
\vspace{0.8em}
{\large Case Brief\par}
\end{center}
\vspace{1em}

\section{Memory Jogger}

In \href{https://www.courtlistener.com/opinion/112404/employment-div-dept-of-human-resources-of-ore-v-smith/}{\textit{Employment Division v. Smith}}, the Supreme Court held that the First Amendment's Free Exercise Clause does not exempt individuals from complying with valid, neutral laws of general applicability, even when those laws incidentally burden central religious practices such as the sacramental use of peyote.

\section{Detailed Case Facts}

Alfred Smith and Galen Black were employed as drug and alcohol rehabilitation counselors by the Douglas County Council on Alcohol and Drug Abuse Prevention and Treatment (ADAPT), a private nonprofit rehabilitation organization in Oregon (\href{https://caselaw.findlaw.com/court/us-supreme-court/485/660.html}{485 U.S. at 662}; \href{https://scholarship.law.uwyo.edu/cgi/viewcontent.cgi?article=1898&context=land_water}{Land \& Water Law Review}). ADAPT maintained a strict drug-free policy under which any employee's use of an illegal nonprescription drug was grounds for immediate termination. Both Smith and Black were members of the Native American Church, and their religious sincerity was never disputed. In 1983 and 1984, each ingested a small quantity of peyote---a hallucinogen derived from the cactus \textit{Lophophora williamsii}---for sacramental purposes during a religious ceremony of the Native American Church (\href{https://supreme.justia.com/cases/federal/us/494/872/}{494 U.S. at 874}). After learning of their peyote use, ADAPT discharged both men (\href{https://tile.loc.gov/storage-services/service/ll/usrep/usrep485/usrep485660/usrep485660.pdf}{485 U.S. at 662--663}).

Smith and Black applied to the Employment Division of the Oregon Department of Human Resources for unemployment compensation. The agency denied their claims under Oregon law disqualifying employees discharged for work-related ``misconduct'' (\href{https://supreme.justia.com/cases/federal/us/494/872/}{494 U.S. at 874}). Although Oregon never criminally prosecuted Smith or Black, Oregon Revised Statutes \S~475.992(4)(a) (1987) made knowing or intentional possession of a Schedule I controlled substance a Class B felony unless prescribed by a medical practitioner, and the Oregon Board of Pharmacy listed peyote on Schedule I under Ore. Admin. Rule 855-80-021(3)(s) (1988) (\href{https://supreme.justia.com/cases/federal/us/494/872/}{494 U.S. at 874}).

Historically and politically, the case arose at the height of the late-1980s ``War on Drugs'' and amidst a long tension between federal accommodation of indigenous religions---reflected in the federal Drug Enforcement Administration's regulatory exemption for Native American Church peyote use, 21 C.F.R. \S~1307.31 (1989), exemptions in 23 states, and Congress's American Indian Religious Freedom Act of 1978, 42 U.S.C. \S~1996 (\href{https://supreme.justia.com/cases/federal/us/494/872/}{494 U.S. at 912 n.5, 920})---and the Rehnquist Court's growing skepticism toward judicially crafted constitutional exemptions under the First Amendment's Free Exercise Clause (applicable to the states via the Fourteenth Amendment under \textit{Cantwell v. Connecticut}, 310 U.S. 296 (1940)).

\section{Procedural History}

\begin{itemize}
  \item \textbf{Administrative Proceedings:} The Oregon Employment Division denied unemployment benefits to Smith and Black for work-related misconduct. After referee hearings reached favorable outcomes for the claimants, the Employment Appeals Board reversed and denied benefits (\href{https://scholarship.law.uwyo.edu/cgi/viewcontent.cgi?article=1898&context=land_water}{Land \& Water Law Review}).
  \item \textbf{Oregon Court of Appeals and Oregon Supreme Court (First Round):} The Oregon Court of Appeals reversed the benefit denials as violating the Free Exercise Clause. The Oregon Supreme Court affirmed in \textit{Smith v. Employment Div., Dept. of Human Resources}, 301 Or. 209, 721 P.2d 445 (1986) (and companion case \textit{Black v. Employment Div.}, 301 Or. 221 (1986)), holding that the potential criminality of peyote use was irrelevant because Oregon's unemployment misconduct rule aimed only to protect the financial integrity of the compensation fund, which could not satisfy strict scrutiny under \textit{Sherbert v. Verner}, 374 U.S. 398 (1963), and \textit{Thomas v. Review Board}, 450 U.S. 707 (1981) (\href{https://supreme.justia.com/cases/federal/us/494/872/}{494 U.S. at 875}).
  \item \textbf{U.S. Supreme Court (\textit{Smith I}):} The Supreme Court granted certiorari, 480 U.S. 916 (1987), and in \href{https://caselaw.findlaw.com/court/us-supreme-court/485/660.html}{\textit{Employment Div., Dept. of Human Resources of Oregon v. Smith}}, 485 U.S. 660 (1988) (\textit{Smith I}), vacated and remanded. The Court held that if Oregon constitutionally prohibited sacramental peyote use through its criminal law, it could impose the lesser burden of denying unemployment benefits; because Oregon's courts had not decided whether state criminal law actually banned sacramental peyote use, the Court remanded for that threshold state-law determination (\href{https://supreme.justia.com/cases/federal/us/494/872/}{494 U.S. at 875--876}).
  \item \textbf{Oregon Supreme Court on Remand:} On remand, the Oregon Supreme Court held that Oregon's controlled-substance statute ``makes no exception for the sacramental use'' of peyote, but concluded that this criminal prohibition violated the federal Free Exercise Clause and reaffirmed its award of unemployment benefits, 307 Or. 68, 763 P.2d 146 (1988) (\href{https://supreme.justia.com/cases/federal/us/494/872/}{494 U.S. at 876}).
  \item \textbf{U.S. Supreme Court (\textit{Smith II}):} The Court again granted certiorari, 489 U.S. 1077 (1989), heard oral argument on November 6, 1989, and on April 17, 1990, reversed the judgment of the Oregon Supreme Court (\href{https://www.courtlistener.com/opinion/112404/employment-div-dept-of-human-resources-of-ore-v-smith/}{494 U.S. at 874, 890}).
\end{itemize}

\section{Judicial Votes}

The Supreme Court reversed by a \textbf{6--3} vote in the judgment, while splitting \textbf{5--4} on the governing Free Exercise standard:

\begin{itemize}
  \item \textbf{Majority (5):} Justice Antonin Scalia delivered the opinion of the Court, joined by Chief Justice William H. Rehnquist and Justices Byron R. White, John Paul Stevens, and Anthony M. Kennedy (\href{https://supreme.justia.com/cases/federal/us/494/872/}{494 U.S. at 874--890}).
  \item \textbf{Concurrence in the Judgment (1; Parts I \& II joined by 3 dissenters):} Justice Sandra Day O'Connor filed an opinion concurring in the judgment (\href{https://supreme.justia.com/cases/federal/us/494/872/}{494 U.S. at 891--907}). Justices William J. Brennan, Jr., Thurgood Marshall, and Harry A. Blackmun joined Parts I and II of Justice O'Connor's opinion (defending strict scrutiny under \textit{Sherbert}) without concurring in the judgment (Part III).
  \item \textbf{Dissent (3):} Justice Harry A. Blackmun filed a dissenting opinion, joined by Justices William J. Brennan, Jr. and Thurgood Marshall (\href{https://supreme.justia.com/cases/federal/us/494/872/}{494 U.S. at 907--921}).
\end{itemize}

\section{Holding}

\textbf{Judgment:} Reversed. Because Oregon's criminal prohibition of peyote use is constitutional under the Free Exercise Clause, Oregon may deny unemployment compensation benefits to employees dismissed for sacramental peyote use.

\textbf{Rule Announced:} The Free Exercise Clause does not relieve an individual of the obligation to comply with a ``valid and neutral law of general applicability'' on the ground that the law proscribes or prescribes conduct that the individual's religion prescribes or proscribes (\href{https://supreme.justia.com/cases/federal/us/494/872/}{494 U.S. at 879}). The compelling-interest balancing test of \textit{Sherbert v. Verner} does not apply to across-the-board, religion-neutral laws, and religious exemptions from such laws are left to the political process rather than mandated by the First Amendment (\href{https://supreme.justia.com/cases/federal/us/494/872/}{494 U.S. at 884--890}).

\textbf{Holding vs. Dicta:} The core holding is that a neutral, generally applicable criminal prohibition (here, Oregon's ban on peyote possession) is valid under the Free Exercise Clause without satisfying strict scrutiny, and therefore supports denying unemployment benefits for violating it. Two important qualifications in Justice Scalia's opinion---carving out (1) ``hybrid'' claims combining Free Exercise with another constitutional right such as free speech or parental rights (\textit{Cantwell}, \textit{Wisconsin v. Yoder}), and (2) schemes with ``a system of individual exemptions'' such as unemployment ``good cause'' rules (\textit{Sherbert}) (\href{https://supreme.justia.com/cases/federal/us/494/872/}{494 U.S. at 881--884})---were explanatory distinctions of prior precedent rather than necessary to decide Smith's claim, yet both became major doctrinal battlegrounds in later cases such as \href{https://www.oyez.org/cases/2020/19-123}{\textit{Fulton v. City of Philadelphia}}, 593 U.S. 522 (2021).

\section{Analysis of Opinions}

\subsection{Majority Opinion (Justice Scalia)}

Justice Scalia began with the text of the Free Exercise Clause, acknowledging that government may never regulate religious beliefs as such or ban physical acts ``only when they are engaged in for religious reasons'' (\href{https://supreme.justia.com/cases/federal/us/494/872/}{494 U.S. at 877}). When a law is neutral and generally applicable, however, an incidental burden on religious conduct is not ``prohibiting the free exercise'' of religion, just as a general tax on businesses does not abridge freedom of the press (\href{https://supreme.justia.com/cases/federal/us/494/872/}{494 U.S. at 878}).

Surveying more than a century of precedent---anchored by \textit{Reynolds v. United States}, 98 U.S. 145 (1879), \textit{Minersville School District v. Gobitis}, 310 U.S. 586 (1940), \textit{Prince v. Massachusetts}, 321 U.S. 158 (1944), \textit{Braunfeld v. Brown}, 366 U.S. 599 (1961), \textit{Gillette v. United States}, 401 U.S. 437 (1971), and \textit{United States v. Lee}, 455 U.S. 252 (1982)---Scalia asserted that the Court had never exempted religiously motivated conduct from a neutral, generally applicable law based on Free Exercise alone (\href{https://supreme.justia.com/cases/federal/us/494/872/}{494 U.S. at 878--881}). Cases that appeared to do so, such as \textit{Cantwell v. Connecticut} and \textit{Wisconsin v. Yoder}, 406 U.S. 205 (1972), were recharacterized as ``hybrid'' cases pairing Free Exercise with free speech, press, or parental rights (\href{https://supreme.justia.com/cases/federal/us/494/872/}{494 U.S. at 881--882}).

Scalia next cabined \textit{Sherbert v. Verner}'s compelling-interest test, observing that the Court had never used it to invalidate any government action outside unemployment compensation and had declined to apply it in \textit{Bowen v. Roy}, 476 U.S. 693 (1986), \textit{Lyng v. Northwest Indian Cemetery Protective Ass'n}, 485 U.S. 439 (1988), \textit{Goldman v. Weinberger}, 475 U.S. 503 (1986), and \textit{O'Lone v. Estate of Shabazz}, 482 U.S. 342 (1987) (\href{https://supreme.justia.com/cases/federal/us/494/872/}{494 U.S. at 883--884}). Whereas unemployment rules invited ``individualized governmental assessment'' of ``good cause,'' an across-the-board criminal prohibition did not. Applying true strict scrutiny to all religious objections would produce a ``constitutional anomaly''---a private right to ignore general laws---and ``court[] anarchy'' in a religiously diverse nation, inviting exemptions from military service, taxes, vaccination, child-neglect, minimum-wage, and antidiscrimination laws (\href{https://supreme.justia.com/cases/federal/us/494/872/}{494 U.S. at 885--889}). Nor could courts limit strict scrutiny to ``central'' religious practices, because judges have no competence to evaluate theological centrality (\href{https://supreme.justia.com/cases/federal/us/494/872/}{494 U.S. at 886--887}). Scalia candidly conceded that leaving religious accommodations to legislatures would place minority faiths at a ``relative disadvantage,'' but declared that ``unavoidable consequence of democratic government'' preferable to making each conscience ``a law unto itself'' (\href{https://supreme.justia.com/cases/federal/us/494/872/}{494 U.S. at 890}).

\subsection{Concurring Opinion in the Judgment (Justice O'Connor; Parts I \& II joined by Justices Brennan, Marshall, and Blackmun)}

Justice O'Connor sharply rebuked the majority for abandoning settled Free Exercise doctrine unnecessarily. In Parts I and II (joined by the three dissenters), she argued that the First Amendment's text makes no distinction between targeted bans and generally applicable laws: a law criminalizing an act of worship manifestly prohibits free exercise regardless of whether it also bars secular conduct (\href{https://supreme.justia.com/cases/federal/us/494/872/}{494 U.S. at 893--894}). She rejected Scalia's reinterpretation of \textit{Cantwell} and \textit{Yoder} as ``hybrid'' cases, noted that \textit{Prince}, \textit{Braunfeld}, \textit{Gillette}, and \textit{Lee} had weighed competing state interests rather than declaring strict scrutiny inapplicable, and distinguished \textit{Roy}, \textit{Lyng}, \textit{Goldman}, and \textit{O'Lone} as specialized internal-government, military, or prison contexts (\href{https://supreme.justia.com/cases/federal/us/494/872/}{494 U.S. at 895--901}). Invoking Justice Jackson's opinion in \textit{West Virginia Board of Education v. Barnette}, 319 U.S. 624 (1943), O'Connor stressed that the very purpose of the Bill of Rights is to withdraw minority liberties from majoritarian politics (\href{https://supreme.justia.com/cases/federal/us/494/872/}{494 U.S. at 902--903}).

Writing for herself alone in Part III, however, O'Connor concluded that Oregon's peyote ban survived strict scrutiny. While acknowledging the severe burden on Native American Church members, she found that Oregon had a compelling interest in preventing physical harm and trafficking associated with Schedule I controlled substances, and that uniform enforcement was essential to accomplish that overriding interest, even though Congress and other states had chosen to grant exemptions (\href{https://supreme.justia.com/cases/federal/us/494/872/}{494 U.S. at 903--907}).

\subsection{Dissenting Opinion (Justice Blackmun, joined by Justices Brennan and Marshall)}

Joining Parts I and II of Justice O'Connor's opinion, Justice Blackmun accused the majority of ``effectuat[ing] a wholesale overturning of settled law'' and overreacting to the nation's drug crisis (\href{https://supreme.justia.com/cases/federal/us/494/872/}{494 U.S. at 908}). Disagreeing with Part III of O'Connor's concurrence, Blackmun argued that strict scrutiny requires comparing interests on the same plane of generality: not Oregon's abstract interest in fighting the ``war on drugs,'' but its narrow interest in refusing a religious exemption for ceremonial peyote use (\href{https://supreme.justia.com/cases/federal/us/494/872/}{494 U.S. at 909--910}).

Blackmun marshaled concrete record evidence to dismantle Oregon's asserted compelling interest: (1) Oregon had never sought to prosecute Smith or Black, rendering its interest merely symbolic; (2) the State offered no evidence of physical harm from sacramental peyote use, which the Native American Church strictly confined to supervised rituals while promoting family responsibility and abstinence from alcohol; (3) illegal trafficking in peyote was virtually nonexistent (DEA seizures totaled only 19.4 pounds nationwide between 1980 and 1987, compared to over 15 million pounds of marijuana); and (4) the federal government and 23 states already maintained peyote exemptions without unleashing a flood of drug-exemption claims (\href{https://supreme.justia.com/cases/federal/us/494/872/}{494 U.S. at 911--918}). Denying an exemption to a faith whose central sacrament is peyote threatened its very survival and mocked federal policy under the American Indian Religious Freedom Act (\href{https://supreme.justia.com/cases/federal/us/494/872/}{494 U.S. at 919--921}).

\section{Examples: Future Applications}

\begin{enumerate}
  \item \textbf{Same-Side (Law Upheld Without Strict Scrutiny) --- Mandatory Childhood Vaccination:} A state enacts a uniform public-school vaccination mandate that allows medical exemptions only where clinically contraindicated, with no religious or philosophical exemptions. Under \textit{Smith}, the mandate is neutral and generally applicable as to religious motivation (assuming medical exemptions are not treated as undermining general applicability), so parents with sincere religious objections have no Free Exercise right to a constitutional exemption and must seek accommodation from the legislature.
  \item \textbf{Same-Side (Law Upheld Without Strict Scrutiny) --- General Zoning Ban on Commercial Structures:} A municipality enforces a neutral zoning ordinance barring all assembly halls---secular and religious alike---from a quiet residential district. Under \textit{Smith}, because the ordinance is not targeted at religion and applies across the board, a congregation cannot claim a Free Exercise Clause exemption merely because the rule incidentally burdens its desire to worship in that zone (though statutory protections such as RLUIPA may independently apply).
  \item \textbf{Opposite-Side (Strict Scrutiny Triggered) --- Targeted Ban on Ritual Animal Slaughter:} A city council enacts ordinances banning the ``ritual'' sacrifice of animals while exempting hunting, pest control, and commercial slaughterhouses, enacted specifically after a Santeria church announces plans to open in the city. Under \textit{Smith}'s neutrality requirement (\href{https://supreme.justia.com/cases/federal/us/494/872/}{494 U.S. at 877--878}), a law targeting religious conduct is neither neutral nor generally applicable and triggers strict scrutiny---the exact scenario later invalidated in \textit{Church of the Lukumi Babalu Aye, Inc. v. City of Hialeah}, 508 U.S. 520 (1993).
  \item \textbf{Opposite-Side (Strict Scrutiny Triggered) --- Discretionary Individualized Exemptions in Contracting:} A city's standard foster-care contract bars agencies from declining to certify same-sex couples, but expressly authorizes a city commissioner to grant exceptions in his ``sole discretion.'' Under \textit{Smith}'s carve-out for systems of individualized exemptions (\href{https://supreme.justia.com/cases/federal/us/494/872/}{494 U.S. at 884}), the city's refusal to extend an exception to a religious agency triggers strict scrutiny---as the Supreme Court held unanimously in \href{https://www.oyez.org/cases/2020/19-123}{\textit{Fulton v. City of Philadelphia}}, 593 U.S. 522 (2021).
  \item \textbf{Fence-Sitter --- Hybrid Claim Combining Free Exercise and Free Speech:} A state public-accommodations statute prohibits all businesses open to the public from refusing service based on sexual orientation, with no discretionary exemptions, and is applied to a religious graphic designer who refuses to create custom wedding websites for same-sex couples. Under \textit{Smith}, the statute is facially neutral and generally applicable, pointing toward upholding the law; yet the designer asserts a ``hybrid'' claim combining Free Exercise with compelled-speech protections (\href{https://supreme.justia.com/cases/federal/us/494/872/}{494 U.S. at 881--882}), leaving lower courts divided over whether \textit{Smith}'s hybrid-rights dictum actually lowers the bar for Free Exercise or adds nothing beyond an independent Free Speech claim.
\end{enumerate}

\section{Critique}

Few modern Supreme Court decisions provoked as swift and broad a bipartisan backlash as \textit{Employment Division v. Smith}. Within three years, a sweeping coalition of civil liberties groups and religious organizations persuaded Congress to pass the Religious Freedom Restoration Act of 1993 (RFRA) by a nearly unanimous vote to restore the \textit{Sherbert} compelling-interest test (\href{https://scholarship.law.gwu.edu/cgi/viewcontent.cgi?article=2806&context=faculty_publications}{GWU Law Scholarship}; \href{https://www.heritage.org/religious-liberty/report/the-religious-freedom-restoration-act-turns-30-the-most-precious-all}{Heritage Foundation}). After the Court struck down RFRA as applied to the states in \textit{City of Boerne v. Flores}, 521 U.S. 507 (1997), Congress responded with the Religious Land Use and Institutionalized Persons Act of 2000 (RLUIPA), and over twenty states enacted their own state-level RFRAs (\href{https://supreme.justia.com/cases/federal/us/494/872/}{Justia Commentary}).

Doctrinally, scholars have identified several major logical weaknesses in Scalia's opinion. First, the ``hybrid rights'' distinction used to salvage \textit{Yoder} and \textit{Cantwell} has been widely criticized as an unprincipled ad hoc invention: if the companion right (such as speech) independently invalidates the law, the Free Exercise Clause does no work; if it does not, combining two losing claims into a winning one lacks coherent justification. Second, \textit{Smith}'s ``individualized exemptions'' carve-out has proven deeply unstable, enabling the modern Court in \href{https://www.scotusblog.com/cases/fulton-v-city-of-philadelphia-pennsylvania/}{\textit{Fulton v. City of Philadelphia}} (2021) and pandemic-era shadow-docket rulings to hollow out ``general applicability'' whenever a law contains even a single secular exception. Third, the parties in \textit{Smith} never briefed or argued whether \textit{Sherbert} should be overruled, making the majority's doctrinal revolution procedurally abrupt.

From an \textbf{originalist and textualist perspective}, \textit{Smith} has fractured conservative jurisprudence. Notably, Scalia's opinion engaged in almost no Founding-era historical analysis of the Free Exercise Clause. Leading originalist scholar Michael McConnell argued that Founding-era colonial charters, state constitutions, and James Madison's writings support a baseline understanding of constitutionally required religious exemptions---an argument later embraced by Justices Alito, Thomas, and Gorsuch in \href{https://www.oyez.org/cases/2020/19-123}{\textit{Fulton}} when calling for \textit{Smith} to be overruled (\href{https://scholarship.law.gwu.edu/cgi/viewcontent.cgi?article=2806&context=faculty_publications}{GWU Law Scholarship}). Conversely, other originalists (such as Philip Hamburger) and judicial-restraint conservatives defend \textit{Smith}'s result, contending that the Founding generation understood ``free exercise'' as protecting against targeted persecution rather than conferring judicially managed exemptions from general civil laws, and warning---as Justice Barrett did in her \textit{Fulton} concurrence---that overruling \textit{Smith} without a clear replacement standard would plunge courts back into subjective balancing.

From a \textbf{progressive perspective}, attitudes toward \textit{Smith} have undergone a striking historical realignment. In 1990, progressives and civil liberties advocates fiercely condemned \textit{Smith} for sacrificing vulnerable minority faiths---especially Native American religious practices---to majoritarian indifference. Over the past two decades, however, as Free Exercise claims have increasingly been invoked by religious institutions and business owners seeking exemptions from LGBTQ+ antidiscrimination laws, contraception mandates, and public-health rules, many progressive scholars have come to view \textit{Smith}'s neutral-and-generally-applicable rule as a vital bulwark protecting equal citizenship and third-party rights against expansive religious exemptions (\href{https://scholarship.law.gwu.edu/cgi/viewcontent.cgi?article=2806&context=faculty_publications}{GWU Law Scholarship}).

\section{Key Quotations}

\begin{quote}
``We have never held that an individual's religious beliefs excuse him from compliance with an otherwise valid law prohibiting conduct that the State is free to regulate.''
--- Justice Scalia, Opinion of the Court (\href{https://supreme.justia.com/cases/federal/us/494/872/}{494 U.S. at 878--879}).
\end{quote}

\begin{quote}
``Precisely because `we are a cosmopolitan nation made up of people of almost every conceivable religious preference,' \textit{Braunfeld v. Brown}, 366 U.S. at 606, and precisely because we value and protect that religious divergence, we cannot afford the luxury of deeming presumptively invalid, as applied to the religious objector, every regulation of conduct that does not protect an interest of the highest order.''
--- Justice Scalia, Opinion of the Court (\href{https://supreme.justia.com/cases/federal/us/494/872/}{494 U.S. at 888}).
\end{quote}

\begin{quote}
``It may fairly be said that leaving accommodation to the political process will place at a relative disadvantage those religious practices that are not widely engaged in; but that unavoidable consequence of democratic government must be preferred to a system in which each conscience is a law unto itself or in which judges weigh the social importance of all laws against the centrality of all religious beliefs.''
--- Justice Scalia, Opinion of the Court (\href{https://supreme.justia.com/cases/federal/us/494/872/}{494 U.S. at 890}).
\end{quote}

\begin{quote}
``There is nothing talismanic about neutral laws of general applicability or general criminal prohibitions, for laws neutral toward religion can coerce a person to violate his religious conscience or intrude upon his religious duties just as effectively as laws aimed at religion.''
--- Justice O'Connor, Concurring in the Judgment (\href{https://supreme.justia.com/cases/federal/us/494/872/}{494 U.S. at 901}).
\end{quote}

\begin{quote}
``I do not believe the Founders thought their dearly bought freedom from religious persecution a `luxury,' but an essential element of liberty --- and they could not have thought religious intolerance `unavoidable,' for they drafted the Religion Clauses precisely in order to avoid that intolerance.''
--- Justice Blackmun, Dissenting (\href{https://supreme.justia.com/cases/federal/us/494/872/}{494 U.S. at 909}).
\end{quote}

\end{document}
